\section{Institutional Interpretation and Empirical Implications}\label{sec:institutions}

The model is written as a theory of manufacturing production under localized climate risk, but its primitives correspond to recognizable policy and business-continuity decisions. The institutional evidence is used to discipline interpretation, not to claim that the exact strategic channel has already been identified empirically.

\subsection{Local climate adaptation and business continuity}

Japan's Ministry of the Environment explicitly assigns local governments a role in regional climate-adaptation planning and, in its 2026 revision of the local adaptation-plan manual, discusses progress management and joint plan formation by multiple local governments \citep{MOE2026}. This supports the model's treatment of place-based public adaptation as a decentralized policy choice.

Firms also make real geographic continuity choices. The Small and Medium Enterprise Agency's business-continuity guidance lists alternative facilities, parallel operation at another site, temporary replacement facilities, inter-firm arrangements, and construction of facilities at another location as possible responses to disruption \citep{SMEA_BCP}. These are direct institutional counterparts to the model's reduced-form backup-readiness choice. The model does not require every firm literally to own a duplicate plant; it requires costly preparations that preserve alternative production capability when a primary site fails.

\subsection{Competition for mobile production}

The Ministry of Economy, Trade and Industry maintains industrial-site information, supports industrial-site development, and facilitates matching between firms and candidate locations \citep{METI2026}. Local governments also explicitly pursue business attraction for employment and regional economic activity. For example, Ehime Prefecture describes enterprise attraction as a policy for employment expansion and regional economic revitalization \citep{EHIME2026}. These facts support the $b$ term in \eqref{eq:government-payoff}: primary production has locally valued benefits that need not be additional national benefits when activity is relocated across jurisdictions.

The evidence does \emph{not} establish that local climate-protection spending is currently chosen for the purpose of attracting firms. That direct causal link is an implication to be tested, not an assumed fact.

\subsection{Empirical predictions}

The theory yields several testable implications, with the scope of each comparative static stated explicitly.
\begin{enumerate}[label=(\roman*)]
\item Improvements in place-based protection should increase primary production exposure to the protected jurisdiction while reducing firms' use of geographically separate contingency capacity, other things equal.
\item On the maintained symmetric interior branch, industries with closer substitutes have lower equilibrium backup readiness because the sole-survivor versus duopoly profit wedge is larger. This is a \emph{level} prediction. The model does not generally predict that product substitutability strengthens the marginal crowd-out effect of public protection. At the canonical co-located symmetric origin, the exact cross-partial is negative (Appendix~\ref{app:verification}), so the readiness reduction induced by an incremental increase in protection is locally \emph{weaker}, not stronger, when products are closer substitutes at this point.
\item Jurisdictions placing a larger value on plant attraction---for example because local employment or immobile-factor rents are more salient---should exhibit a larger decentralized protection incentive relative to a coordinated benchmark.
\item Evaluation of regional resilience projects based only on engineering reductions in local failure probability may overstate net industrial resilience benefits if firms respond by reducing costly alternative production arrangements.
\end{enumerate}
These predictions distinguish the mechanism from a pure safe-development story in which only the exposure of households or physical capital changes after public protection. Prediction (ii) deliberately separates the level effect of product substitutability from the interaction between substitutability and the marginal policy response; the latter is not assigned a global sign by the model.
